\chapter*{Úvod}
\phantomsection
\addcontentsline{toc}{chapter}{Úvod}


Chytrá domácnost je dnes tvořena zařízeními od různých výrobců, která se bez
společného standardu těžko propojují. Tento problém řeší protokol
Matter~\cite{csa-matter}, který definuje jednotný aplikační model zařízení
nad IP vrstvou. Přenos dat po rádiovém rozhraní IEEE~802.15.4 zajišťuje
nízkopříkonová síť Thread, postavená na IPv6~\cite{thread-group}.

Tato bakalářská práce se věnuje návrhu a realizaci tří prvků chytré domácnosti,
které komunikují protokolem Matter přenášeným přes síť Thread. Jde o čidlo
teploty, chytré světlo a bateriový vypínač. Všechna zařízení jsou postavena
na čipech Espressif ESP32-C6, které disponují rádiem IEEE~802.15.4 a podporou
protokolu Matter~\cite{espressif-esp32c6-ds}. Pro vývoj byl zvolen modul
Seeed Studio XIAO ESP32-C6~\cite{seeed-xiao-wiki}.